% !PRINTED 158
\weberpage{158}

%<<CONT
forms that can be reduced to one another only with
 difficulty. First we obtain from the definition (1),
if we now denote the parameter by $a$:

\weq{6}{\begin{aligned}
&\Pi(u+v,a)-\Pi(u,a)-\Pi(v,a)\\
&\qquad = \dfrac{1}{2}\log\dfrac{\Theta(u+v-a)}{\Theta(u+v+a)}\,\dfrac{\Theta(u+a)\,\Theta(v+a)}{\Theta(u-a)\,\Theta(v-a)},
\end{aligned}}

and it remains only to represent the expression standing under the
logarithm by means of elliptic functions.

This is first done easily, as at the passage of the Fundamenta already
mentioned, by means of the formula [§ 43, (11)]:

\[
\Theta^2(0)\,\Theta(u+v)\,\Theta(u-v) = \Theta^2(u)\,\Theta^2(v)\,[1-\varkappa^2\,\mathrm{sn}^2u\,\mathrm{sn}^2v],
\]

if one puts therein for $u,v$ the values $u\pm a$, $v\pm a$, then $\pm a$,
$u+v\pm a$.

One thus obtains:

\[
\begin{aligned}
\Theta^2(0)\,\Theta(u+v\pm 2a)\,\Theta(u-v) &= \Theta^2(u\pm a)\,\Theta^2(v\pm a)\,[1-\varkappa^2\,\mathrm{sn}^2(u\pm a)\,\mathrm{sn}^2(v\pm a)],\\
\Theta^2(0)\,\Theta(u+v\pm 2a)\,\Theta(u+v) &= \Theta^2(a)\,\Theta^2(u+v\pm a)\,[1-\varkappa^2\,\mathrm{sn}^2a\,\mathrm{sn}^2(u+v\pm a)],
\end{aligned}
\]

from which it follows:

\weq{7}{\dfrac{\Theta(u+v-a)\,\Theta(u+a)\,\Theta(v+a)}{\Theta(u+v+a)\,\Theta(u-a)\,\Theta(v-a)} = \sqrt{\dfrac{[1-\varkappa^2\,\mathrm{sn}^2(u-a)\,\mathrm{sn}^2(v-a)]\,[1-\varkappa^2\,\mathrm{sn}^2a\,\mathrm{sn}^2(u+v+a)]}{[1-\varkappa^2\,\mathrm{sn}^2(u+a)\,\mathrm{sn}^2(v+a)]\,[1-\varkappa^2\,\mathrm{sn}^2a\,\mathrm{sn}^2(u+v-a)]}}.}

A second expression is obtained from the addition formula
(13), § 22, which can be put in the form:

\[
\Theta(0)\,\Theta(u\pm a)\,\Theta(v\pm a)\,\Theta(u+v) = \Theta(u)\,\Theta(v)\,\Theta(a)\,\Theta(u+v\pm a)\,[1\pm\varkappa^2\,\mathrm{sn}\,a\,\mathrm{sn}\,u\,\mathrm{sn}\,v\,\mathrm{sn}(u+v\pm a)],
\]

from which by division:

\weq{8}{\dfrac{\Theta(u+v-a)\,\Theta(u+a)\,\Theta(v+a)}{\Theta(u+v+a)\,\Theta(u-a)\,\Theta(v-a)} = \dfrac{1+\varkappa^2\,\mathrm{sn}\,a\,\mathrm{sn}\,u\,\mathrm{sn}\,v\,\mathrm{sn}(u+v+a)}{1-\varkappa^2\,\mathrm{sn}\,a\,\mathrm{sn}\,u\,\mathrm{sn}\,v\,\mathrm{sn}(u+v-a)}.}

\gesperrt{Jacobi} gives yet a third expression for the
$\Theta$-quotients in (6):

\weq{9}{\dfrac{\left\{1-\varkappa^2\,\mathrm{sn}^2\!\left(\dfrac{u-v}{2}\right)\mathrm{sn}^2\!\left(\dfrac{u+v}{2}+a\right)\right\}\left\{1-\varkappa^2\,\mathrm{sn}^2\!\left(\dfrac{u+v}{2}\right)\mathrm{sn}^2\!\left(\dfrac{u+v}{2}-a\right)\right\}}{\left\{1-\varkappa^2\,\mathrm{sn}^2\!\left(\dfrac{u-v}{2}\right)\mathrm{sn}^2\!\left(\dfrac{u+v}{2}-a\right)\right\}\left\{1-\varkappa^2\,\mathrm{sn}^2\!\left(\dfrac{u+v}{2}\right)\mathrm{sn}^2\!\left(\dfrac{u+v}{2}+a\right)\right\}}.}
